% 数学ノート集・共通テンプレート v2（積分論2026年度第1部の版面を参考）
% A5判。章内で定義・定理・例・注意等の番号を共通化。
% コンパイル： platex -interaction=nonstopmode math_note_template_v2.tex (2回)
%              dvipdfmx math_note_template_v2.dvi
\documentclass[a5paper,10pt,onecolumn]{jsarticle}
\usepackage[dvipdfmx,top=22mm,bottom=20mm,inner=19mm,outer=18mm,headheight=13pt,headsep=6mm]{geometry}
\usepackage{amsmath,amssymb,amsthm,mathtools}
\usepackage{enumitem,fancyhdr,titlesec}
\usepackage[dvipdfmx,bookmarks=true,colorlinks=false]{hyperref}
\usepackage{pxjahyper}

% ----- 差し替える文書情報 -----
\newcommand{\NoteTitle}{算術関数・Dirichlet畳み込みとゼータ関数}
\newcommand{\NoteSeries}{数学ノート集}

% ----- 文字・組版 -----
% 参考PDFの質素な明朝＋Computer Modernの印象を引き継ぎながら，
% A5判で窮屈にならないよう本文の行間を少し確保する．
\linespread{1.13}
\setlength{\parindent}{1zw}
\setlength{\parskip}{0pt}
\setlength{\abovedisplayskip}{6pt plus 2pt minus 2pt}
\setlength{\belowdisplayskip}{6pt plus 2pt minus 2pt}
\setlength{\abovedisplayshortskip}{3pt plus 2pt minus 1pt}
\setlength{\belowdisplayshortskip}{4pt plus 2pt minus 1pt}
\setlist[enumerate]{topsep=3pt,itemsep=1pt,parsep=0pt,leftmargin=2.4zw}
\titleformat{\section}{\normalfont\normalsize\bfseries}{\thesection}{0.8zw}{}
\titlespacing*{\section}{0pt}{2.0ex plus .4ex}{.9ex}
\titleformat{\subsection}{\normalfont\normalsize\bfseries}{\thesubsection}{0.7zw}{}
\titlespacing*{\subsection}{0pt}{1.5ex plus .3ex}{.6ex}
\raggedbottom

% ----- 柱とページ番号：ページ上部、罫線なし -----
\fancypagestyle{mathnote}{%
  \fancyhf{}%
  \fancyhead[L]{\footnotesize \NoteSeries\ （\NoteTitle）}%
  \fancyhead[R]{\footnotesize \thepage}%
  \renewcommand{\headrulewidth}{0pt}%
  \renewcommand{\footrulewidth}{0pt}%
}
\pagestyle{mathnote}

% ----- 節内通し番号の定義・定理群 -----
\newtheoremstyle{mathnotetheorem}
 {4.5pt}{4.5pt}{\normalfont}{}%
 {\bfseries}{}{0.65zw}%
 {\thmname{#1}\thmnumber{ #2}\thmnote{（#3）}}
\theoremstyle{mathnotetheorem}
\newtheorem{theorem}{定理}[section]
\newtheorem{definition}[theorem]{定義}
\newtheorem{lemma}[theorem]{補題}
\newtheorem{proposition}[theorem]{命題}
\newtheorem{corollary}[theorem]{系}
\newtheorem{example}[theorem]{例}
\newtheorem{remark}[theorem]{注意}
\renewcommand{\proofname}{証明}
\renewcommand{\qedsymbol}{$\square$}
\makeatletter
\renewenvironment{proof}[1][\proofname]{%
  \par\pushQED{\qed}\normalfont\topsep4\p@\@plus2\p@\relax
  \trivlist\item[\hskip\labelsep\bfseries #1]\ignorespaces
}{\popQED\endtrivlist\@endpefalse}
\makeatother

% ----- 使うノートで適宜拡張する記号 -----
\newcommand{\R}{\mathbb{R}}
\newcommand{\norm}[1]{\lVert #1\rVert}
\newcommand{\abs}[1]{\lvert #1\rvert}

\begin{document}
\section*{読了後にできること}
本稿を読むための前提は，算術関数，約数和，Dirichlet畳み込み．
本稿の具体的な到達課題は，Dirichlet反転を使って算術関数を復元し，具体的な値を計算する．


\section{算術関数と乗法性}

本稿では自然数は $1$ から始まるものとする．
複素数値算術関数とは，写像
$f:\mathbb N_{\ge1}\to\mathbb C$ をいう．

\begin{definition}
算術関数 $f$ が\emph{乗法的}であるとは
$f(1)=1$ かつ
\[
 f(mn)=f(m)f(n)\qquad(\gcd(m,n)=1)
\]
を満たすことをいう．
また全ての自然数 $m,n$ について上の等式が成立するとき
\emph{完全乗法的}という．
\end{definition}

\begin{definition}
$n$ の正の約数の個数を $d(n)$，
正の約数の総和を $\sigma(n)$，
$n$ と互いに素な $1\le k\le n$ の個数を
$\varphi(n)$ と書く．
\end{definition}

\begin{proposition}[素数冪上の値]\label{prop:basic-values}
素数 $p$ と整数 $a\ge1$ に対して
\[
 d(p^a)=a+1,\qquad
 \sigma(p^a)=\frac{p^{a+1}-1}{p-1},\qquad
 \varphi(p^a)=p^a-p^{a-1}.
\]
さらに $d,\sigma,\varphi$ は乗法的である．
\end{proposition}
\begin{proof}
$p^a$ の正の約数は
$1,p,\ldots,p^a$ なので前二式を得る．
$1\le k\le p^a$ のうち $p$ と互いに素でないものは
$p$ の倍数であり，その個数は $p^{a-1}$ だから
$\varphi(p^a)=p^a-p^{a-1}$ となる．
$\gcd(m,n)=1$ なら任意の $mn$ の約数は一意に
$m$ の約数と $n$ の約数の積に分解される．
従って $d,\sigma$ は乗法的である．
また中国剰余定理によって
$(\mathbb Z/mn\mathbb Z)^\times$ と
$(\mathbb Z/m\mathbb Z)^\times
 \times(\mathbb Z/n\mathbb Z)^\times$
が同型となる．
両群の位数を比較すれば
$\varphi(mn)=\varphi(m)\varphi(n)$ を得る．
\end{proof}

\section{Dirichlet畳み込み}

\begin{definition}
算術関数 $f,g$ の\emph{Dirichlet畳み込み}を
\[
 (f*g)(n)=\sum_{d\mid n}f(d)g(n/d)
\]
と定める．
また
$\mathbf1(n)=1$，
$\varepsilon(1)=1$，
$\varepsilon(n)=0$（$n>1$）と置く．
\end{definition}

\begin{proposition}[畳み込みの代数構造]\label{prop:convolution-ring}
算術関数全体は点ごとの加法と
Dirichlet畳み込みについて可換環をなす．
その乗法単位元は $\varepsilon$ である．
\end{proposition}
\begin{proof}
畳み込みの交換律は約数 $d$ と $n/d$ の
置換によって従う．
結合律は
\[
 ((f*g)*h)(n)=\sum_{abc=n}f(a)g(b)h(c)
 =(f*(g*h))(n)
\]
から従う．この和は有限であるため，
和の順序の変更に問題はない．
分配律は有限和の線形性による．
また
\[
 (f*\varepsilon)(n)
 =\sum_{d\mid n}f(d)\varepsilon(n/d)=f(n)
\]
だから $\varepsilon$ は単位元である．
\end{proof}

\begin{lemma}[畳み込みの逆元]\label{lem:unit}
算術関数 $f$ がDirichlet畳み込みについて
逆元を持つための必要十分条件は $f(1)\ne0$ である．
\end{lemma}
\begin{proof}
$f*g=\varepsilon$ なら
$f(1)g(1)=1$ なので $f(1)\ne0$ である．
逆に $f(1)\ne0$ とし，
$g(1)=1/f(1)$ と定める．
$n>1$ では
\[
 0=(f*g)(n)
 =f(1)g(n)+\sum_{\substack{d\mid n\\d>1}}
 f(d)g(n/d)
\]
となることを要求する．
$1\le n/d<n$ だから，
右辺の和は帰納的に既に定まった $g(k)$ だけを含む．
従って
\[
 g(n)=-\frac1{f(1)}
 \sum_{\substack{d\mid n\\d>1}}f(d)g(n/d)
\]
によって $g(n)$ を一意に定義できる．
この帰納法は全ての $n$ で
$f*g=\varepsilon$ を保証する．
\end{proof}

\begin{proposition}[乗法性の保存]\label{prop:multiplicative}
$f,g$ が乗法的なら $f*g$ も乗法的である．
\end{proposition}
\begin{proof}
$\gcd(m,n)=1$ とする．
$mn$ の約数は一意に $a\mid m$，$b\mid n$ に対して
$ab$ と書けるので
\begin{align*}
 (f*g)(mn)
 &=\sum_{a\mid m}\sum_{b\mid n}
 f(ab)g(mn/(ab))\\
 &=\sum_{a\mid m}\sum_{b\mid n}
 f(a)f(b)g(m/a)g(n/b)\\
 &=(f*g)(m)(f*g)(n).
\end{align*}
また $(f*g)(1)=f(1)g(1)=1$ である．
\end{proof}

\section{M\"obius反転}

\begin{definition}
\emph{M\"obius関数} $\mu$ を
\[
 \mu(1)=1,\qquad
 \mu(n)=\begin{cases}
 (-1)^r&\text{$n$ が相異なる $r$ 個の素数の積},\\
 0&\text{$n$ が平方因子を持つ}
 \end{cases}
\]
と定める．
\end{definition}

\begin{lemma}[M\"obius関数の基本恒等式]\label{lem:mu}
\[
 \sum_{d\mid n}\mu(d)
 =\begin{cases}1&(n=1),\\0&(n>1)\end{cases}.
\]
すなわち $\mathbf1*\mu=\varepsilon$ である．
\end{lemma}
\begin{proof}
$n=1$ では左辺は $\mu(1)=1$ である．
$n>1$ では相異なる素因数を
$p_1,\ldots,p_r$ とすると $r\ge1$ である．
$\mu(d)\ne0$ となる約数 $d$ は
$p_1,\ldots,p_r$ の部分集合の積に限られる．
従って
\[
 \sum_{d\mid n}\mu(d)
 =\sum_{k=0}^r(-1)^k\binom rk
 =(1-1)^r=0.
\]
\end{proof}

\begin{theorem}[M\"obius反転公式]\label{thm:mobius}
算術関数 $f,g$ について
\[
 g(n)=\sum_{d\mid n}f(d)
\]
が全ての $n$ で成立するための必要十分条件は
\[
 f(n)=\sum_{d\mid n}\mu(d)g(n/d)
\]
が全ての $n$ で成立することである．
\end{theorem}
\begin{proof}
最初の等式は $g=\mathbf1*f$ と書ける．
両辺に $\mu$ を畳み込むと
\[
 \mu*g=(\mu*\mathbf1)*f
 =\varepsilon*f=f
\]
となる．
逆に $f=\mu*g$ なら
$\mathbf1*f=(\mathbf1*\mu)*g
 =\varepsilon*g=g$ である．
\end{proof}

\begin{corollary}[Euler関数の約数恒等式]\label{cor:phi}
任意の自然数 $n$ について
\[
 \sum_{d\mid n}\varphi(d)=n,\qquad
 \varphi(n)=n\sum_{d\mid n}\frac{\mu(d)}d
 =n\prod_{p\mid n}\left(1-\frac1p\right)
\]
が成り立つ．
\end{corollary}
\begin{proof}
$1\le k\le n$ を $\gcd(k,n)=d$ によって分類する．
$d\mid n$ を固定すると
$k=du$，$1\le u\le n/d$，
$\gcd(u,n/d)=1$ となる $u$ は
$\varphi(n/d)$ 個である．
よって $n=\sum_{d\mid n}\varphi(n/d)$ を得る．
M\"obius反転から
\[
 \varphi(n)=\sum_{d\mid n}\mu(d)\frac nd
 =n\sum_{d\mid n}\frac{\mu(d)}d.
\]
最後の和は，$\mu(d)$ が非零となる約数を
素因数の部分集合で指定することにより
$\prod_{p\mid n}(1-1/p)$ に等しい．
\end{proof}

\section{Dirichlet級数}

\begin{definition}
算術関数 $f$ に対して，
複素数 $s$ を変数とする級数
\[
 D_f(s)=\sum_{n=1}^\infty\frac{f(n)}{n^s}
\]
を\emph{Dirichlet級数}という．
ここで $n^{-s}=e^{-s\log n}$ とする．
\end{definition}

\begin{lemma}[絶対収束領域と正則性]\label{lem:dirichlet-holo}
ある $\sigma_0\in\mathbb R$ について
$\sum_n|f(n)|n^{-\sigma_0}<\infty$ とする．
このとき $\operatorname{Re}s>\sigma_0$ で
$D_f(s)$ は正則であり，
任意の非負整数 $k$ に対して
\[
 D_f^{(k)}(s)=\sum_{n=1}^\infty
 \frac{(-\log n)^kf(n)}{n^s}
\]
が成り立つ．
\end{lemma}
\begin{proof}
コンパクト集合 $C\subset\{\operatorname{Re}s>\sigma_0\}$
を固定する．ある $\delta>0$ が存在して
$\operatorname{Re}s\ge\sigma_0+\delta$ が
全ての $s\in C$ で成り立つ．
任意の $k$ について
$(\log n)^kn^{-\delta}$ は $n\ge1$ で有界だから，
\[
 \left|\frac{(-\log n)^kf(n)}{n^s}\right|
 \le C_{k,\delta}|f(n)|n^{-\sigma_0}.
\]
右辺の級数は収束するので，
WeierstrassのM判定法によって
元の級数と全ての微分級数は
$C$ 上で一様収束する．
正則関数の局所一様収束定理と項別微分定理を
繰り返し適用すると結論が従う．
\end{proof}

\begin{proposition}[畳み込みは積に対応する]\label{prop:series-product}
$f,g$ のDirichlet級数がある $s$ で絶対収束するとき，
\[
 D_{f*g}(s)=D_f(s)D_g(s)
\]
が成り立つ．
\end{proposition}
\begin{proof}
絶対収束の仮定から
\[
 \sum_{a,b\ge1}
 \left|\frac{f(a)g(b)}{(ab)^s}\right|
 =\left(\sum_a\frac{|f(a)|}{a^{\operatorname{Re}s}}\right)
  \left(\sum_b\frac{|g(b)|}{b^{\operatorname{Re}s}}\right)
 <\infty.
\]
従って二重級数の和の順序を変更でき，
\begin{align*}
 D_f(s)D_g(s)
 &=\sum_{a,b\ge1}\frac{f(a)g(b)}{(ab)^s}\\
 &=\sum_{n\ge1}\frac1{n^s}
   \sum_{ab=n}f(a)g(b)
 =D_{f*g}(s).
\end{align*}
\end{proof}

\section{Riemannゼータ関数とEuler積}

\begin{definition}
$\operatorname{Re}s>1$ に対して
\[
 \zeta(s)=\sum_{n=1}^\infty n^{-s}
\]
と定める．
これをRiemannゼータ関数という．
\end{definition}

\begin{theorem}[Euler積]\label{thm:euler-product}
$\operatorname{Re}s>1$ において
\[
 \zeta(s)=\prod_{p\ {\rm prime}}(1-p^{-s})^{-1}
\]
が成立し，右辺は有限個の素数による部分積の極限として
局所一様収束する．
特にこの半平面で $\zeta(s)\ne0$ である．
\end{theorem}
\begin{proof}
$\sigma=\operatorname{Re}s>1$ とする．
有限集合の素数 $P$ について
\[
 \prod_{p\in P}(1-p^{-s})^{-1}
 =\prod_{p\in P}\sum_{k=0}^\infty p^{-ks}
 =\sum_{\substack{n\ge1\\
   \text{全ての素因数が }P\text{ に属する}}}n^{-s}.
\]
各幾何級数は絶対収束しており，
有限個の積と和の入れ替えは正当化される．
整数の素因数分解の一意性を用いると，
全ての素数を順に追加した極限は
$\sum_{n\ge1}n^{-s}=\zeta(s)$ に等しい．
なぜなら除外された項の絶対値の和は
$\sum_{n>N}n^{-\sigma}$ によって上から評価でき，
$P$ が $N$ 以下の全ての素数を含む場合には
その和は $N\to\infty$ で零となるからである．
$\sigma\ge1+\delta$ となるコンパクト集合上では
同じ評価により局所一様収束する．

零点を持たないことを示す．
$|p^{-s}|\le p^{-\sigma}<1$ なので，
\[
 \prod_p(1-p^{-s})
\]
は $\sum_pp^{-\sigma}<\infty$ から
非零の極限に収束する．
実際十分大きい $p$ では
$|\log(1-p^{-s})|\le2p^{-\sigma}$ であり，
対数級数が絶対収束するからである．
従ってEuler積の極限も非零であり，
$\zeta(s)\ne0$ となる．
\end{proof}

\begin{corollary}[M\"obius関数のDirichlet級数]
\label{cor:mu-dirichlet}
$\operatorname{Re}s>1$ では
\[
 \sum_{n=1}^\infty\frac{\mu(n)}{n^s}
 =\frac1{\zeta(s)}
\]
が成り立つ．
\end{corollary}
\begin{proof}
$|\mu(n)|\le1$ なので左辺は絶対収束する．
$\mathbf1*\mu=\varepsilon$ と
命題\ref{prop:series-product}から
\[
 \zeta(s)\sum_n\mu(n)n^{-s}
 =\sum_n\varepsilon(n)n^{-s}=1.
\]
Euler積によって $\zeta(s)\ne0$ なので結論が従う．
\end{proof}

\begin{proposition}[対数微分]\label{prop:log-derivative}
$\operatorname{Re}s>1$ では
\[
 -\frac{\zeta'(s)}{\zeta(s)}
 =\sum_p\sum_{k=1}^\infty
 \frac{\log p}{p^{ks}}
 =\sum_{n=1}^\infty\frac{\Lambda(n)}{n^s}
\]
が成立する．
ここで $\Lambda(n)=\log p$（$n=p^k,\ k\ge1$），
それ以外は $\Lambda(n)=0$ とする．
\end{proposition}
\begin{proof}
$\operatorname{Re}s\ge1+\delta$ のコンパクト集合上で
\[
 \sum_p\sum_{k\ge1}
 \frac{\log p}{p^{k\operatorname{Re}s}}
 \le\sum_{n\ge2}\frac{\log n}{n^{1+\delta}}
      \frac1{1-2^{-(1+\delta)}}
 <\infty
\]
と評価できる．
従ってEuler積に対して
絶対収束する対数級数の項別微分を適用してよい．
$\log\zeta(s)=-\sum_p\log(1-p^{-s})$ と
$\frac d{ds}p^{-s}=-(\log p)p^{-s}$
から
\[
 \frac{\zeta'(s)}{\zeta(s)}
 =-\sum_p\frac{(\log p)p^{-s}}{1-p^{-s}}
 =-\sum_p\sum_{k\ge1}(\log p)p^{-ks}.
\]
最後に素数冪 $p^k$ に対応する項をまとめれば
$\Lambda(n)$ による表示式となる．
\end{proof}

\section{素数の逆数和}

\begin{theorem}[Eulerの定理]\label{thm:prime-reciprocals}
素数全体にわたる級数
\[
 \sum_{p\ {\rm prime}}\frac1p
\]
は発散する．
\end{theorem}
\begin{proof}
$\sigma>1$ とする．Euler積の対数を取ると
\[
 \log\zeta(\sigma)
 =\sum_p\sum_{k=1}^\infty
 \frac{p^{-k\sigma}}k
 =\sum_pp^{-\sigma}+R(\sigma),
\]
ただし
\[
 0\le R(\sigma)
 =\sum_p\sum_{k=2}^\infty\frac{p^{-k\sigma}}k
 \le\sum_{n=2}^\infty
     \frac{n^{-2\sigma}}{1-n^{-\sigma}}
 \le\frac1{1-2^{-1}}
     \sum_{n=2}^\infty n^{-2}
\]
である．従って $R(\sigma)$ は
$\sigma>1$ で一様に有界である．
一方，単調減少関数の積分比較により
\[
 \zeta(\sigma)=\sum_{n=1}^\infty n^{-\sigma}
 \ge\int_1^\infty t^{-\sigma}\,dt
 =\frac1{\sigma-1}
\]
だから $\sigma\downarrow1$ で
$\log\zeta(\sigma)\to\infty$ である．
ゆえに $\sum_pp^{-\sigma}\to\infty$ となる．
もし $\sum_pp^{-1}$ が収束するなら
$p^{-\sigma}\le p^{-1}$ によって
$\sum_pp^{-\sigma}$ は一様有界になるはずで矛盾する．
従って結論を得る．
\end{proof}

\section{約数関数の平均}

\begin{lemma}[調和数列]\label{lem:harmonic}
\[
 H_N=\sum_{n=1}^N\frac1n
 =\log N+\gamma+O(N^{-1})
\]
となる実定数 $\gamma$ が存在する．
\end{lemma}
\begin{proof}
$a_N=H_N-\log N$ と置く．
\[
 a_{N+1}-a_N
 =\frac1{N+1}-\log\left(1+\frac1N\right)
\]
である．
積分表示
$\log(1+t)=\int_0^t(1+u)^{-1}du$ と
$1-u\le(1+u)^{-1}\le1$（$0\le u\le1$）
から $|a_{N+1}-a_N|\le C/N^2$ が従う．
従って $(a_N)$ はCauchy列であり，
ある実数 $\gamma$ へ収束する．
また
\[
 |a_N-\gamma|
 \le\sum_{k=N}^\infty|a_{k+1}-a_k|
 \le C\sum_{k=N}^\infty k^{-2}=O(N^{-1}).
\]
よって表示式が成立する．
\end{proof}

\begin{theorem}[Dirichletの約数問題の初等評価]
\label{thm:divisor-average}
$x\ge2$ とし，
$D(x)=\sum_{n\le x}d(n)$ と置く．
このとき
\[
 D(x)=x\log x+(2\gamma-1)x+O(\sqrt x)
\]
が成り立つ．
\end{theorem}
\begin{proof}
$d(n)=\sum_{ab=n}1$ なので
\[
 D(x)=\#\{(a,b)\in\mathbb N^2:ab\le x\}.
\]
$m=\lfloor\sqrt x\rfloor$ と置く．
$ab\le x$ なら $a\le m$ または $b\le m$ である．
二つの集合を数え，重なりの $m^2$ 個を引くと
\[
 D(x)=2\sum_{a=1}^m\left\lfloor\frac xa\right\rfloor-m^2.
\]
各項について $\lfloor x/a\rfloor=x/a+O(1)$ なので
\[
 D(x)=2xH_m-m^2+O(m).
\]
補題\ref{lem:harmonic}から
$H_m=\log m+\gamma+O(1/m)$ である．
$m=\sqrt x+O(1)$ だから
\[
 \log m=\tfrac12\log x+O(x^{-1/2}),\qquad
 m^2=x+O(\sqrt x).
\]
従って
\[
 D(x)=x\log x+2\gamma x-x+O(\sqrt x)
\]
となり，結論を得る．
\end{proof}

\begin{remark}
本稿で得たEuler積と対数微分は
解析的整数論の出発点である．
Dirichletの算術級数定理では指標と $L$ 関数が，
素数定理では $\zeta(s)$ の境界上の零点の情報が
さらに必要となる．
本稿ではこれらの深い定理を証明したことにはしない．
\end{remark}

\section{到達課題の完全解答}
\noindent\textbf{Möbius反転の計算．}
\begin{proposition}
$\mathbf1(n)=1$，$\varepsilon(1)=1$，
$\varepsilon(n)=0$（$n>1$）とすると
$\mathbf1*\mu=\varepsilon$ である．
\end{proposition}
\begin{proof}
$(\mathbf1*\mu)(n)=\sum_{d\mid n}\mu(d)$．
$n=1$ の値は1である．
$n>1$ の相異なる素因子が $r\ge1$ 個あれば，
平方因子を含まない約数の寄与は
$\sum_{j=0}^r\binom rj(-1)^j=(1-1)^r=0$ である．
\end{proof}
\begin{proposition}
$\sum_{d\mid n}\varphi(d)=n$ から
$\varphi(n)=n\prod_{p\mid n}(1-1/p)$ が従う．
\end{proposition}
\begin{proof}
$\varphi*\mathbf1=\mathrm{id}$ と表せるので
$\varphi=\mathrm{id}*\mu$．従って
$\varphi(n)=\sum_{d\mid n}d\mu(n/d)
=n\sum_{d\mid n}\mu(d)/d$．
平方因子を含まない約数に分けると有限積を得る．
\end{proof}
\noindent\textbf{到達確認．}
$\varphi(60)$ を計算せよ．
\par\noindent\textbf{解答．}
$60=2^2\cdot3\cdot5$ より
$\varphi(60)=60(1-1/2)(1-1/3)(1-1/5)=16$．
\end{document}